\cotefile{94}
\scanpage{14}\mspage{- II -}
(corps de fonctions algébriques) et un rapport sur la Géométrie algébrique

A ce propos le rédacteur signale que , s'il n'a pas eu lui-même de trop grandes difficultés à donner un traitement presque \textquotedbl{}\unsure{gleichberechtigt}\textquotedbl{} des $p$-adiques et séries formelles , son collègue global a dû séparer nombres et fonctions , et ceci pour des raisons qui semblent \deleted{\unsure{valables}} sérieuses Ceci ne peut qu'apporter de l'eau au moulin de ceux qui veulent traiter les corps de fonctions au moyen de courbes , produits de courbes et jacobiennes . Cependant le rédacteur \unsure{demande} qu'on ne s'engage pas à la légère à renoncer à l'analogie entre fonctions et nombres .

Le rédacteur signale les détails suivants :

a) La prop.3 du \S~1 pourrait être améliorée par prolongement de la valeur absolue d'une algèbre normée réelle à sa complexifiée .

b) Les résultats sur les extensions non ramifiées (\S~2,fin du n\textsuperscript{o}3) méritent des énoncés en forme .

c) Le n\textsuperscript{o} des représentants multiplicatifs (\S~2,n\textsuperscript{o}5) est inutile pour la suite Celui de l'exponentielle \deleted{\illegible{1}} (\S~2,n\textsuperscript{o}6) l'est presque , n'étaient l'utilisation de la finitude de $(K:K^{n})$ dans le th. d'existence du corps de classes local (\S~6,n\textsuperscript{o}7) , et le besoin de la valeur de cet indice qu'a eu le collègue global. Par contre l'exponentielle est inutile pour l'inégalité $(K:N(L))\quad L:K$ ($K$ localement compact , $L$ galoisienne) , qui est fondamentale pour le corps de classes : on peut démontrer celle-ci , soit par une astuce krasnérienne (cf.\S~6 n\textsuperscript{o}1) ; soit comme conséquence de la Führerdiskriminantenformel (\S~6,n\textsuperscript{o}8) ; ce n'est pas nettement plus compliqué que la méthode Chevalley .

d) Le n\textsuperscript{o} du groupe de décomposition est mal fichu (\S~3,n\textsuperscript{o}1) . Il faudrait des énoncés en forme au \S~3,n\textsuperscript{o}3 .

e) Le rédacteur est affolé par les amas de cocycles \deleted{\unsure{sur}} du corps de classes local , et ne comprend rien à tous ces calculs \quad Il a l'impression que la méthode des algèbres simples rend les choses plus claires . Mais il sera ravi si le Haut Commissariat à la Topologie algébrique prend les \deleted{\unsure{\S\S|de}} n\textsuperscript{os} 2,3 et 4 du \S~6 à son compte et y met de l'ordre .

\begin{center}\texttt{-{}-{}-{}-{}-{}-{}-{}-{}-{}-{}-{}-{}-{}-{}-{}-{}-{}-{}-}\end{center}
% archive sheet number: 13
